% appendix.tex -- Supplementary Material

\section{Sequence Encoding, PCA Projection, and Multiplicity Conditioning}\label{app:pca}

Stochastic attention (SA) operates on continuous vectors on the unit sphere,
while protein sequences are discrete objects over a 20-letter amino acid
alphabet. Bridging the two requires three transformations (one-hot encoding,
PCA dimensionality reduction, and unit-norm projection), and the reverse path
(inverse PCA followed by argmax decoding) recovers discrete sequences from
continuous samples. Here we describe the full encoding--decoding pipeline,
derive the multiplicity-weighted conditioning formulas used in the main text,
and discuss how the multiplicity vector $\mathbf{r}$ interacts with the PCA
representation.

Given a cleaned alignment of $K$ sequences, each of length $L$ amino acids
over the standard 20-letter alphabet
$\mathcal{A} = \{\text{A}, \text{R}, \text{N}, \ldots, \text{V}\}$, we
encoded each sequence as a binary vector in $\mathbb{R}^{20L}$ by
concatenating per-position indicator vectors. For sequence $k$ ($k = 1,
\ldots, K$) with amino acid $a_\ell^{(k)}$ at alignment position $\ell$
($\ell = 1, \ldots, L$):
\begin{equation}
  \mathbf{x}_k = \bigl[\mathbf{e}_{a_1^{(k)}},\; \mathbf{e}_{a_2^{(k)}},\;
    \ldots,\; \mathbf{e}_{a_L^{(k)}}\bigr]^\top \in \{0,1\}^{20L},
\end{equation}
where $\mathbf{e}_a \in \{0,1\}^{20}$ is the standard basis vector for amino
acid $a$ (i.e., a one-hot indicator with a 1 in the position corresponding to
$a$ and 0 elsewhere). The full one-hot dimensionality is
$d_{\mathrm{full}} = 20L$, which ranged from 520 ($\omega$-conotoxin,
$L{=}26$) to 1{,}060 (Kunitz, $L{=}53$) across the four families. This space
is over-dimensioned relative to the number of stored patterns $K$
($d_{\mathrm{full}}/K$ ranged from 7.0 to 10.7), which creates a problem for
the modern Hopfield energy~\cite{ramsauerHopfieldNetworksAll2021}: the
similarity score $e_k = \mathbf{m}_k^\top\boldsymbol{\xi}$ between a stored
pattern $\mathbf{m}_k$ and a query state $\boldsymbol{\xi}$ has variance
$1/d$ for a random query on the unit sphere $\mathbb{S}^{d-1}$, so when $d$
is large all scores concentrate near zero and the inverse temperature $\beta$
must be set extremely high to differentiate among patterns, making Langevin
sampling inefficient. PCA removes this redundancy by projecting onto the
$d \ll d_{\mathrm{full}}$ directions of actual variation in the alignment.

We computed PCA from the one-hot encoded alignment by first centering the data:
$\bar{\mathbf{x}} = K^{-1}\sum_{k=1}^K \mathbf{x}_k$ (the family's
position-specific amino acid frequency profile),
$\tilde{\mathbf{X}} = \mathbf{X} - \bar{\mathbf{x}}\mathbf{1}_K^\top$,
and then computing the economy SVD
$\tilde{\mathbf{X}} = \mathbf{U}\boldsymbol{\Sigma}\mathbf{V}^\top$, where
$\mathbf{U} \in \mathbb{R}^{d_{\mathrm{full}} \times r}$ contains the
principal directions of variation,
$\boldsymbol{\Sigma} = \mathrm{diag}(\sigma_1, \ldots, \sigma_r)$ the
singular values in decreasing order, and
$r = \mathrm{rank}(\tilde{\mathbf{X}}) \leq \min(d_{\mathrm{full}}, K) - 1$.
We retained the smallest dimension $d$ such that
$\sum_{j=1}^d \sigma_j^2 / \sum_{j=1}^r \sigma_j^2 \geq 0.95$ (at least
95\% of total variance). Across the four families, $d$ ranged from 34
($\omega$-conotoxin) to 186 (WW), yielding compression ratios of
$3.3\times$ to $21\times$ (Table~\ref{tab:cross-family}). Each sequence was
then projected as
$\mathbf{z}_k = \mathbf{W}_d^\top(\mathbf{x}_k - \bar{\mathbf{x}})
\in \mathbb{R}^d$, where
$\mathbf{W}_d = \mathbf{U}_{:,1:d}$ is the projection matrix, and
normalized to unit $L_2$ norm:
$\mathbf{m}_k = \mathbf{z}_k / \|\mathbf{z}_k\|_2 \in \mathbb{S}^{d-1}$.
This normalization ensures that similarity scores lie in $[-1, 1]$ and
prevents patterns with larger PCA norms from dominating the softmax attention.
The columns $\mathbf{m}_1, \ldots, \mathbf{m}_K$ form the memory matrix
$\hat{\mathbf{X}} \in \mathbb{R}^{d \times K}$ used throughout the paper.
The Langevin sampler (Eq.~\ref{eq:ula}) produces continuous states
$\boldsymbol{\xi} \in \mathbb{R}^d$, which are decoded back to amino acid
sequences by inverse PCA reconstruction
$\hat{\mathbf{x}} = \bar{\mathbf{x}} + \mathbf{W}_d\boldsymbol{\xi}
\in \mathbb{R}^{d_{\mathrm{full}}}$ followed by per-position argmax:
$a_\ell = \arg\max_{a \in \mathcal{A}}\; \hat{x}_{20(\ell-1) + a}$ for
$\ell = 1, \ldots, L$. This decoding is deterministic and produced valid
amino acids at 100\% of positions across all families and conditions.

The multiplicity-weighted Hopfield energy (Eq.~\ref{eq:weighted-energy})
assigns a weight $r_k > 0$ to each stored pattern $\mathbf{m}_k$. For the
binary functional split used throughout this paper, we set $r_k = \rho$ for
the $K_{\mathrm{des}}$ designated patterns and $r_k = 1$ for the
$K_{\mathrm{bg}} = K - K_{\mathrm{des}}$ background patterns, where
$\rho \geq 1$ is the multiplicity ratio. The weighted score function
(Proposition~\ref{prop:score}) adds $\log\rho$ to the pre-softmax logits of
designated patterns and $\log 1 = 0$ to background patterns. To derive the
effect of $\rho$ on the attention distribution, consider the idealized case
where the query $\boldsymbol{\xi}$ has equal similarity to all stored
patterns, i.e., $\mathbf{m}_k^\top\boldsymbol{\xi} = c$ for all $k$. The
softmax attention weight on pattern $k$ simplifies to
$a_k = r_k / \sum_{j} r_j$, and the total attention weight on all designated
patterns is:
\begin{equation}\label{eq:feff-def}
  f_{\mathrm{eff}}(\rho)
  \;=\; \sum_{k \in \mathrm{des}} a_k
  \;=\; \frac{K_{\mathrm{des}}\,\rho}{K_{\mathrm{des}}\,\rho + K_{\mathrm{bg}}}.
\end{equation}
At $\rho = 1$, $f_{\mathrm{eff}} = K_{\mathrm{des}} / K$ (the natural
proportion); as $\rho \to \infty$, $f_{\mathrm{eff}} \to 1$. To find the
multiplicity ratio needed for a target effective fraction $f \in (0, 1)$, we
invert Eq.~\eqref{eq:feff-def} by multiplying both sides by the denominator
and solving for $\rho$:
\begin{align}
  f\,(K_{\mathrm{des}}\,\rho + K_{\mathrm{bg}}) &= K_{\mathrm{des}}\,\rho,
  \notag\\
  f\,K_{\mathrm{bg}} &= K_{\mathrm{des}}\,\rho\,(1 - f), \notag\\
  \rho(f) &= \frac{f\,K_{\mathrm{bg}}}{K_{\mathrm{des}}\,(1-f)}.
  \label{eq:rho-inverse-derived}
\end{align}
This is Eq.~\eqref{eq:rho-inverse} in the main text. The formula is exact
under the equal-similarity assumption; in practice, we confirmed empirically
that the mean attention weight on designated patterns tracked
$f_{\mathrm{eff}}(\rho)$ with $< 0.3\%$ deviation across all $\rho$ values
and all families (Table~\ref{tab:per-family-rho}), validating the
approximation as operationally exact. The effective number of patterns under
multiplicity weighting,
$K_{\mathrm{eff}}(\mathbf{r}) = (\sum_k r_k)^2 / \sum_k r_k^2
= (K_{\mathrm{des}}\,\rho + K_{\mathrm{bg}})^2 /
(K_{\mathrm{des}}\,\rho^2 + K_{\mathrm{bg}})$, equals $K$ when $\rho = 1$
and converges to $K_{\mathrm{des}}$ as $\rho \to \infty$; this quantity
governs the shift in the phase transition temperature $\beta^{*}$.

A key property of multiplicity-weighted conditioning is that $\mathbf{r}$
does \emph{not} alter the PCA encoding. The memory matrix
$\hat{\mathbf{X}}$ is constructed once from the full family alignment and
remains fixed regardless of $\rho$; the multiplicity weights enter only
through the softmax logits. This separation has three consequences.
First, the PCA subspace is independent of $\rho$: the principal directions
$\mathbf{W}_d$, the centering vector $\bar{\mathbf{x}}$, and the retained
dimension $d$ are all determined by the alignment alone, so the calibration
gap $\Delta = f_{\mathrm{eff}} - f_{\mathrm{obs}}$ is identical for all
$\rho$ values on a given family. Second, $\rho$ controls attention, not
geometry: increasing $\rho$ adds $\log\rho$ to the designated logits, shifting
attention weight toward designated patterns, but this shift must still pass
through the PCA reconstruction and argmax decoding to produce discrete
residues, which is where the calibration gap arises. Third, $\rho$ shifts the
phase transition: the logit biases pre-concentrate the attention distribution
at $\beta = 0$, reducing the weighted entropy to
$H_{\mathbf{r}}(0) = \log K_{\mathrm{eff}}(\mathbf{r}) < \log K$, so
reaching the retrieval-phase inflection requires a higher inverse temperature
$\beta^{*}(\rho) \geq \beta^{*}(1)$. We observed this empirically:
$\beta^{*}$ increased from 4.4 ($\rho=1$, $K_{\mathrm{eff}} = 99$) to 9.3
($\rho=1{,}000$, $K_{\mathrm{eff}} = 32.1$) on the Kunitz domain
(Fig.~\ref{fig:phase-transition}). By contrast, hard curation (where the
memory matrix is built from the designated subset only) \emph{does} change
the PCA: the centering vector, principal directions, and retained dimension
are all computed from $K_{\mathrm{des}}$ sequences, capturing
subset-specific variation that may be orthogonal to the leading components
of the full-family PCA. This is why hard curation achieved 100\% phenotype
transfer regardless of the Fisher separation index $S$, while multiplicity
weighting was limited by the full-family PCA geometry.

\section{Supplementary Tables and Figures}\label{app:tables-figures}

Tables~\ref{tab:kunitz-rho}--\ref{tab:forkhead-rho} report selected values from
the $\rho$ sweep on each of the five Pfam benchmark families.
$f_{\mathrm{eff}}$ is the effective designated fraction computed from the
multiplicity vector; $f_{\mathrm{obs}}$ is the hard-decoded marker fraction
averaged over 930 generated sequences; $\bar{a}_{\mathrm{des}}$ is the mean
softmax attention weight on designated patterns; and $D$ is the mean pairwise
diversity. Table~\ref{tab:beta-sweep} reports phenotype fraction and
diversity as a function of $\beta/\beta^{*}$ for three fixed $\rho$ values
on the Kunitz domain. Figure~\ref{fig:plddt-tmscore-scatter} shows the joint
distribution of ESMFold pLDDT and TM-score for all 250 predicted Kunitz
structures: SA-generated sequences cluster tightly in the high-pLDDT,
high-TM-score region while HMM-emitted sequences scatter broadly across
lower values. The entropy curves for the full $\rho$ sweep
($\rho \in \{1, 2, 5, 10, 20, 50, 100, 200, 500, 1000\}$) on the Kunitz
family produce a monotonically rightward-shifting family of sigmoid curves
with $\beta^{*}(\rho)$ increasing from 4.35 ($\rho=1$) to 9.26
($\rho=1000$), consistent with $\beta^{*} \propto \log K_{\mathrm{eff}}$
predicted by the mean-field argument in Section~\ref{sec:theory}
(Fig.~\ref{fig:phase-transition}).

\begin{table}[ht]
  \centering
  \caption{Kunitz domain (PF00014) $\rho$ sweep. $K=99$, $K_{\mathrm{des}}=32$.}
  \label{tab:kunitz-rho}
  \small
  \begin{tabular}{ccccc}
    \toprule
    $\rho$ & $f_{\mathrm{eff}}$ & $f_{\mathrm{obs}}$ & $\bar{a}_{\mathrm{des}}$ & $D$ \\
    \midrule
    1    & 0.323 & 0.406 & 0.324 & 0.563 \\
    5    & 0.705 & 0.481 & 0.712 & 0.556 \\
    10   & 0.827 & 0.539 & 0.832 & 0.547 \\
    50   & 0.962 & 0.598 & 0.963 & 0.526 \\
    100  & 0.980 & 0.615 & 0.981 & 0.520 \\
    200  & 0.990 & 0.623 & 0.991 & 0.519 \\
    500  & 0.996 & 0.631 & 0.997 & 0.515 \\
    1000 & 0.998 & 0.634 & 0.998 & 0.514 \\
    \bottomrule
  \end{tabular}
\end{table}

\begin{table}[ht]
  \centering
  \caption{SH3 domain (PF00018) $\rho$ sweep. $K=55$, $K_{\mathrm{des}}=33$.}
  \label{tab:sh3-rho}
  \small
  \begin{tabular}{ccccc}
    \toprule
    $\rho$ & $f_{\mathrm{eff}}$ & $f_{\mathrm{obs}}$ & $\bar{a}_{\mathrm{des}}$ & $D$ \\
    \midrule
    1    & 0.600 & 0.852 & 0.614 & 0.580 \\
    5    & 0.882 & 0.897 & 0.890 & 0.575 \\
    10   & 0.938 & 0.906 & 0.941 & 0.575 \\
    50   & 0.983 & 0.965 & 0.983 & 0.567 \\
    100  & 0.991 & 0.977 & 0.991 & 0.559 \\
    200  & 0.996 & 0.984 & 0.996 & 0.554 \\
    500  & 0.999 & 0.989 & 0.999 & 0.513 \\
    1000 & 0.999 & 0.992 & 0.999 & 0.510 \\
    \bottomrule
  \end{tabular}
\end{table}

\begin{table}[ht]
  \centering
  \caption{WW domain (PF00397) $\rho$ sweep. $K=420$, $K_{\mathrm{des}}=69$.}
  \label{tab:ww-rho}
  \small
  \begin{tabular}{ccccc}
    \toprule
    $\rho$ & $f_{\mathrm{eff}}$ & $f_{\mathrm{obs}}$ & $\bar{a}_{\mathrm{des}}$ & $D$ \\
    \midrule
    1    & 0.164 & 0.131 & 0.152 & 0.690 \\
    5    & 0.496 & 0.184 & 0.487 & 0.689 \\
    10   & 0.663 & 0.226 & 0.661 & 0.676 \\
    50   & 0.901 & 0.277 & 0.901 & 0.657 \\
    100  & 0.947 & 0.303 & 0.947 & 0.649 \\
    200  & 0.973 & 0.327 & 0.974 & 0.643 \\
    500  & 0.990 & 0.360 & 0.991 & 0.642 \\
    1000 & 0.995 & 0.371 & 0.995 & 0.642 \\
    \bottomrule
  \end{tabular}
\end{table}

\begin{table}[ht]
  \centering
  \caption{Homeobox domain (PF00046) $\rho$ sweep. $K=136$, $K_{\mathrm{des}}=102$.}
  \label{tab:homeobox-rho}
  \small
  \begin{tabular}{ccccc}
    \toprule
    $\rho$ & $f_{\mathrm{eff}}$ & $f_{\mathrm{obs}}$ & $\bar{a}_{\mathrm{des}}$ & $D$ \\
    \midrule
    1    & 0.750 & 0.935 & 0.758 & 0.525 \\
    5    & 0.938 & 0.954 & 0.939 & 0.521 \\
    10   & 0.968 & 0.955 & 0.969 & 0.522 \\
    50   & 0.993 & 0.957 & 0.994 & 0.521 \\
    100  & 0.997 & 0.957 & 0.997 & 0.521 \\
    200  & 0.998 & 0.957 & 0.998 & 0.521 \\
    500  & 0.999 & 0.957 & 0.999 & 0.521 \\
    1000 & 1.000 & 0.957 & 1.000 & 0.521 \\
    \bottomrule
  \end{tabular}
\end{table}

\begin{table}[ht]
  \centering
  \caption{Forkhead domain (PF00250) $\rho$ sweep. $K=246$, $K_{\mathrm{des}}=122$.}
  \label{tab:forkhead-rho}
  \small
  \begin{tabular}{ccccc}
    \toprule
    $\rho$ & $f_{\mathrm{eff}}$ & $f_{\mathrm{obs}}$ & $\bar{a}_{\mathrm{des}}$ & $D$ \\
    \midrule
    1    & 0.496 & 0.556 & 0.490 & 0.480 \\
    5    & 0.831 & 0.610 & 0.826 & 0.464 \\
    10   & 0.908 & 0.638 & 0.904 & 0.447 \\
    50   & 0.980 & 0.670 & 0.979 & 0.431 \\
    100  & 0.990 & 0.710 & 0.988 & 0.417 \\
    200  & 0.995 & 0.711 & 0.994 & 0.416 \\
    500  & 0.998 & 0.730 & 0.998 & 0.399 \\
    1000 & 0.999 & 0.730 & 0.999 & 0.398 \\
    \bottomrule
  \end{tabular}
\end{table}

\begin{table}[ht]
  \centering
  \caption{$\beta$ sweep on Kunitz domain ($K=99$, $K_{\mathrm{des}}=32$).}
  \label{tab:beta-sweep}
  \small
  \begin{tabular}{cccc}
    \toprule
    $\rho$ & $\beta/\beta^{*}$ & $f_{\mathrm{obs}}$ & $D$ \\
    \midrule
    10  & 0.50 & 0.471 & 0.581 \\
    10  & 1.00 & 0.539 & 0.547 \\
    10  & 1.50 & 0.573 & 0.525 \\
    10  & 2.00 & 0.597 & 0.508 \\
    10  & 3.00 & 0.613 & 0.484 \\
    \midrule
    50  & 0.50 & 0.506 & 0.573 \\
    50  & 1.00 & 0.582 & 0.537 \\
    50  & 1.50 & 0.631 & 0.514 \\
    50  & 2.00 & 0.653 & 0.499 \\
    50  & 3.00 & 0.700 & 0.469 \\
    \midrule
    200 & 0.50 & 0.535 & 0.564 \\
    200 & 1.00 & 0.613 & 0.527 \\
    200 & 1.50 & 0.648 & 0.503 \\
    200 & 2.00 & 0.677 & 0.488 \\
    200 & 3.00 & 0.706 & 0.463 \\
  \bottomrule
  \end{tabular}
\end{table}

\section*{Docking-statistics note}
For the $\omega$-conotoxin docking validation (Table~\ref{tab:docking-validation}),
group-level comparisons (SA strong-conditioned, SA full-conditioned, and
natural $\omega$-conotoxin controls) were evaluated using exact two-sided
permutation tests on mean complex-level scores (iPTM, pTM, confidence, and
interface-pLDDT), with sample sizes $n_{\mathrm{SA,strong}}=10$,
$n_{\mathrm{SA,full}}=10$, and $n_{\mathrm{control}}=5$. No pairwise
contrast reached significance at $\alpha=0.05$ (all $p > 0.05$; smallest observed
$p=0.0726$ for confidence, SA full vs controls).

\begin{table}[ht]
  \centering
  \caption{Pairwise exact permutation p-values for $\omega$-conotoxin docking comparisons (all $H_0$: equal means).}
  \label{tab:docking-permutation-pvalues}
  \begin{tabular}{lccc}
    \toprule
    \textbf{Metric} & \textbf{SA strong vs SA full} & \textbf{SA strong vs controls} & \textbf{SA full vs controls} \\
    \midrule
    iPTM & 1.0000 & 0.7502 & 0.4752 \\
    pTM & 0.7873 & 0.1285 & 0.1092 \\
    Confidence & 0.6560 & 0.1245 & 0.0726 \\
    Interface pLDDT & 0.6556 & 0.2238 & 0.2934 \\
    \bottomrule
  \end{tabular}
\end{table}
